\documentclass{article}
\usepackage[utf8]{inputenc}
\usepackage[T1]{fontenc}
\usepackage[french]{babel}
\usepackage{booktabs}
\usepackage{siunitx}
\usepackage{geometry}
\geometry{margin=2cm}

\sisetup{
    round-mode = places,
    round-precision = 1,
    group-separator = {\,},
}

\begin{document}

\section*{Dynamic conic bundle vs sans bundle --- RLT dualis\'e}

Comparaison du temps de calcul total par donn\'ee entre le moteur
\texttt{dynamic conic bundle} (\texttt{engine=conicbundle\_native},
\texttt{dualize=[RLT]}) et le m\^eme mod\`ele r\'esolu sans bundle (MOSEK
classique, RLT comme contrainte dure).

\bigskip
\noindent\textbf{Attention} : la colonne ``sans bundle'' pour UntargetedSDP
provient du fichier \texttt{ECML-SDPu-FP}, dont la configuration diff\`ere
l\'eg\`erement de celle utilis\'ee pour le run dynamic bundle
(\texttt{use\_active\_neurons: false}, et 2 familles de coupes en moins dans
\texttt{cuts}) -- la comparaison n'est donc pas strictement \`a
iso-configuration pour ce mod\`ele. Pour TargetedSDP, le run de r\'ef\'erence
(\texttt{SDPt-RLT=10}) utilise \texttt{RLT\_prop=0.1} (s\'election heuristique)
au lieu d'une dualisation dynamique par pool.

\begin{table}[h]
\centering
\caption{Donn\'ees NON certifi\'ees (sans RLT) -- data\_index $\in \{7, 23, 31, 45, 59\}$}
\begin{tabular}{l S S}
\toprule
{Mod\`ele} & {Temps dynamic (s)} & {Temps sans bundle (s)} \\
\midrule
TargetedSDP   & 876.2  & 2015.6 \\
UntargetedSDP & 194.5  & 4060.9 \\
\bottomrule
\end{tabular}
\end{table}

\begin{table}[h]
\centering
\caption{Donn\'ees certifi\'ees (sans RLT) -- data\_index $\in \{6, 64, 69, 84, 89\}$}
\begin{tabular}{l S S}
\toprule
{Mod\`ele} & {Temps dynamic (s)} & {Temps sans bundle (s)} \\
\midrule
TargetedSDP   & 481.7 & 1541.9 \\
UntargetedSDP & 376.7 & 3429.0 \\
\bottomrule
\end{tabular}
\\[0.3em]
\footnotesize{UntargetedSDP sans bundle : $n=4$ (data\_index=89 manquante dans \texttt{ECML-SDPu-FP}).}
\end{table}

\end{document}
